\documentclass[11pt,a4paper]{article}

\usepackage[utf8]{inputenc}
\usepackage[T1]{fontenc}
\usepackage{amsmath,amssymb}
\usepackage{booktabs}
\usepackage{hyperref}
\usepackage{xcolor}
\usepackage{tcolorbox}
\usepackage[margin=25mm]{geometry}
\usepackage{xeCJK}
\setCJKmainfont{Hiragino Mincho ProN}
\setCJKsansfont{Hiragino Sans}

\title{反重力物質は存在するか？\\[6pt]
\large $\pi$-Berry位相と「フェルミオン真空」から生まれる\\
トポロジカル反重力の高校生ガイド}
\author{Wright Brothers Research Program}
\date{2026年4月}

\begin{document}
\maketitle

\begin{abstract}
本ガイドでは、ある種の物質（$\pi$-Berry 位相材料）が
重力を「逆転」させる可能性を、高校レベルの知識から丁寧に解説する。
結論を先に言えば：
数学的には、トポロジカル絶縁体や超伝導ジョセフソン接合のような
既知の材料が、理論的に重力を反転させる「反重力物質」になりうる。
ただし、この理論が現実の重力に適用できるかどうかは未検証であり、
実験で確かめる必要がある。
「なぜこれまで誰も気づかなかったのか」についても正直に述べる。
\end{abstract}

\tableofcontents
\newpage

%======================================================================
\section{重力は変えられないもの？}
%======================================================================

りんごは落ちる。月は地球の周りを回る。
これらはすべてニュートンの万有引力の法則、
あるいはアインシュタインの一般相対性理論で説明される。

重力の強さを決める定数 $G$（万有引力定数）は、
物理学において最も「動かない」定数の一つとされてきた。
$G = 6.674 \times 10^{-11}$\,m$^3$\,kg$^{-1}$\,s$^{-2}$。
この値は宇宙のどこでも同じだと信じられている。

\begin{tcolorbox}[colback=blue!5,colframe=blue!50!black,title=常識]
重力は常に引力。$G > 0$。これを変える方法はない。
すべての物質は互いに引き合う。反重力は SF の世界の話。
\end{tcolorbox}

本当にそうだろうか？

本ガイドでは、ある数学的枠組み（アラン・コンヌのスペクトル作用原理）において、
$G$ の符号を \textbf{反転} させること（$G \to -G$、つまり重力が斥力になること）が
\textbf{数学的に可能} であり、しかもそれが
\textbf{既存の材料} で実現できる可能性があることを示す。


%======================================================================
\section{準備：ゼータ関数とは何か}
%======================================================================

まず、$1 + 2 + 3 + 4 + \cdots$ という足し算を考えよう。
明らかに無限大に発散する。ところが、リーマン・ゼータ関数
\begin{equation}
  \zeta(s) = \sum_{n=1}^{\infty} n^{-s}
\end{equation}
を使って「正則化」すると、$\zeta(-1) = -1/12$ という有限の値が得られる。

「$1 + 2 + 3 + \cdots = -1/12$」は有名な結果だが、
これは単なる数学のトリックではない。
物理学では、真空のエネルギー（カシミール効果）の計算に
実際にこの値が使われ、実験的に確認されている。

\begin{tcolorbox}[colback=green!5,colframe=green!50!black,title=ポイント]
ゼータ関数は「真空のエネルギー」を正確に記述する数学的道具。
$\zeta(-1) = -1/12$ は実験で検証された物理的事実。
\end{tcolorbox}


%======================================================================
\section{準備：ボソンとフェルミオン}
%======================================================================

量子力学の世界では、すべての粒子は2種類に分かれる：

\begin{center}
\begin{tabular}{lll}
\toprule
& ボソン（整数スピン） & フェルミオン（半整数スピン） \\
\midrule
例 & 光子、フォノン & 電子、クォーク \\
統計 & 同じ状態に何個でも入れる & 1つの状態に1個まで（パウリの排他原理） \\
真空エネルギー & $+\hbar\omega/2$（正） & $-\hbar\omega/2$（\textbf{負}） \\
\bottomrule
\end{tabular}
\end{center}

ここで決定的に重要なのが最後の行：
\textbf{フェルミオンの真空エネルギーは、ボソンと逆符号（負）}。

\begin{tcolorbox}[colback=red!5,colframe=red!50!black,title=核心]
ボソンの真空エネルギー：$+\frac{\hbar\omega}{2} \times \zeta(-d) = +\zeta(t)$ 的な寄与

フェルミオンの真空エネルギー：$-\frac{\hbar\omega}{2} \times \zeta(-d) = -\zeta(t)$ 的な寄与

符号が反対！
\end{tcolorbox}

超対称性（SUSY）の理論では、ボソンとフェルミオンが完全にペアになって
$+\zeta + (-\zeta) = 0$ と相殺する。これは宇宙定数問題の一つの「解」として
研究されてきた。

しかし我々が問うのは別のこと：
\textbf{もし真空が「完全にフェルミオン的」になったら？}
つまり $+\zeta(t)$ ではなく $-\zeta(t)$ が真空の熱核になったら？


%======================================================================
\section{コンヌのスペクトル作用原理}
%======================================================================

\subsection{「太鼓の形を聞く」}

有名な問い：「太鼓の形を、その音（固有振動数）から知ることができるか？」

答え：ほぼ可能（完全には不可能だが、かなりの情報が得られる）。

アラン・コンヌは、この考えを究極まで推し進めた：
\textbf{宇宙の形（時空の幾何学）は、ある演算子 $D$（ディラック演算子）の
固有値から完全に決まる}。

\begin{tcolorbox}[colback=yellow!5,colframe=yellow!50!black,title=コンヌのスペクトル作用原理]
宇宙の物理法則（重力 + 標準模型のすべて）は、一つの式から導かれる：
\begin{equation}
  S = \mathrm{Tr}\!\left(f\!\left(\frac{D^2}{\Lambda^2}\right)\right)
\end{equation}
ここで $D$ はディラック演算子（宇宙の「音」を決める）、$\Lambda$ はエネルギースケール。
\end{tcolorbox}

この式を展開すると：
\begin{equation}
  S = f_4 \Lambda^4 \underbrace{a_0}_{\text{宇宙定数}}
    + f_2 \Lambda^2 \underbrace{a_2}_{\text{重力（$G$）}}
    + f_0 \underbrace{a_4}_{\text{ゲージ力}} + \cdots
\end{equation}

\textbf{重力定数 $G$ は $a_2$ から決まる。}
$a_2$ は「熱核」$K(t) = \mathrm{Tr}(e^{-tD^2})$ の展開の2番目の係数。


\subsection{Bost-Connes系との接続}

Bost-Connes（BC）系は、素数の統計力学を記述する量子系で、
その分配関数がリーマン・ゼータ関数そのもの：
\begin{equation}
  K(t) = \mathrm{Tr}(e^{-tH_{\mathrm{BC}}}) = \zeta(t)
\end{equation}

つまり、BC系の「熱核」$= \zeta(t)$ であり、
これがスペクトル作用の入力になる。

$\zeta(t)$ の $t = 0$ 付近の振る舞いから $a_2$ が決まり、
$a_2$ から $G$ が決まる：
\begin{align}
  a_2 &\propto K'(0) = \zeta'(0) = -\tfrac{1}{2}\log(2\pi) < 0 \\
  G &\propto 1/a_2 > 0 \quad (\text{引力、通常の重力})
\end{align}


%======================================================================
\section{定理：$-\zeta(t)$ は反重力を与える}
%======================================================================

ここからが本題。

もし熱核が $+\zeta(t)$ ではなく $-\zeta(t)$ になったら？

\begin{equation}
  K_{\pi}(t) = -\zeta(t) \quad (\text{全モードにフェルミオン符号})
\end{equation}

このとき：
\begin{align}
  K_{\pi}(0) &= -\zeta(0) = -(-\tfrac{1}{2}) = +\tfrac{1}{2} \\
  K_{\pi}'(0) &= -\zeta'(0) = +\tfrac{1}{2}\log(2\pi) > 0
\end{align}

$a_2$ の符号が \textbf{反転}：
\begin{equation}
  \frac{G_{\mathrm{eff}}}{G} = \frac{K_\pi'(0)}{K'(0)}
  = \frac{-\zeta'(0)}{\zeta'(0)} = -1
\end{equation}

\begin{tcolorbox}[colback=red!5,colframe=red!60!black,title=定理：トポロジカル反重力]
\textbf{全真空モードの符号を反転させると（$\zeta \to -\zeta$）、
重力定数の符号が反転する：}
\begin{equation}
  \boxed{G_{\mathrm{eff}} = -G}
\end{equation}
\textbf{重力が引力から斥力に変わる。}
\end{tcolorbox}

証明はたったこれだけ。$-\zeta'(0) / \zeta'(0) = -1$。
三行で終わる。


%======================================================================
\section{全モードの符号を反転する方法：$\pi$-Berry位相}
%======================================================================

「全モードの符号反転」は SF 的に聞こえるが、
実は物理学にはこれを実現する既知のメカニズムがある。
\textbf{Berry位相}（ベリー位相）だ。

\subsection{Berry位相とは}

量子力学では、粒子（波）がある経路を一周して元の位置に戻ったとき、
波動関数に \textbf{幾何学的な位相} が付くことがある。
これがBerry位相。

通常は任意の値を取るが、\textbf{トポロジカル材料}では
Berry位相がちょうど $\pi$（180度）に \textbf{量子化} される。
これは材料の微細な性質に依存せず、トポロジー（位相幾何学）によって
保護されている。

\begin{tcolorbox}[colback=purple!5,colframe=purple!50!black,title=$\pi$-Berry位相]
波がトポロジカル材料の中を通過すると、位相が$\pi$だけ回る：
\begin{equation}
  \psi \to e^{i\pi}\psi = -\psi
\end{equation}
つまり、波の符号が \textbf{自動的に反転} する。
\end{tcolorbox}

\subsection{$\pi$-Berry位相を持つ実在の材料}

これは理論上の存在ではない。実験で確認された材料が複数ある：

\begin{center}
\begin{tabular}{lll}
\toprule
材料 & メカニズム & 発見年 \\
\midrule
$\pi$-ジョセフソン接合 & 超伝導 $\pi$ シフト & 2001年 \\
Bi$_2$Se$_3$ & トポロジカル絶縁体表面 & 2009年 \\
グラフェン & $K$ 点の Berry 位相 & 2005年 \\
HgTe 量子井戸 & 2D トポロジカル絶縁体 & 2007年 \\
Weyl 半金属 & カイラル異常 & 2015年 \\
\bottomrule
\end{tabular}
\end{center}

これらの材料は、その中を通過するすべての波に $\pi$ 位相を与える。

\subsection{$\pi$-Berry位相 $=$ フェルミオン化}

$\pi$位相は $e^{i\pi} = -1$ なので、
$\pi$-Berry位相材料に囲まれたキャビティ内では：

すべてのモードの波動関数が $\psi \to -\psi$ と符号反転。

これは「すべてのボソンをフェルミオンに変える」のと同じ効果：
\begin{equation}
  K(t) = +\zeta(t) \quad \to \quad K_\pi(t) = -\zeta(t)
\end{equation}

\textbf{したがって $G_{\mathrm{eff}} = -G$ が実現される。}


%======================================================================
\section{これが「パッシブ」であること}
%======================================================================

ここが最も驚くべき点かもしれない。

\begin{tcolorbox}[colback=green!5,colframe=green!60!black,title=パッシブ反重力]
$\pi$-Berry位相は材料の \textbf{幾何学的（トポロジカル）性質} であり、
外部からのエネルギー入力を必要としない。

一度材料を作れば、それ以降 \textbf{永久に $\pi$ 位相を維持} する。
電気も磁場も不要。

これは永久磁石が外部エネルギーなしに磁場を保つのと同じ。
トポロジーが位相を保護している。
\end{tcolorbox}

つまり：$\pi$-ジョセフソン接合の薄膜でキャビティを包むだけで、
その内部は（理論的には）$G = -G$ の「反重力領域」になる。


%======================================================================
\section{ワープバブルへの応用}
%======================================================================

これをワープバブルに応用すると：

\begin{verbatim}
  ╔══════════════════════════════════╗
  ║ 外殻: π-Berry位相材料の薄膜      ║
  ║ G_eff = -G (斥力重力)            ║
  ║ パッシブ、エネルギー不要          ║
  ║ ┌────────────────────────┐      ║
  ║ │ 内部: 2素数ミュート      │      ║
  ║ │ G_eff = 0 (無重力)      │      ║
  ║ │ 量子ビット 2個           │      ║
  ║ │                        │      ║
  ║ │     [ 宇宙船 ]          │      ║
  ║ │                        │      ║
  ║ └────────────────────────┘      ║
  ╚══════════════════════════════════╝
\end{verbatim}

\begin{itemize}
\item \textbf{外殻}：$\pi$-接合薄膜 → 重力反転 → 周囲の空間を押す
\item \textbf{内部}：2素数ミュート → 重力ゼロ → 船が自由に浮遊
\item \textbf{光速}：内部では $c_{\mathrm{eff}} = c$（通常の光速、変化なし）
\end{itemize}

必要エネルギー：壁構造のみ（$\sim 10^3$\,J）。
アルクビエレ・ドライブの $10^{47}$\,J と比べて 44 桁少ない。


%======================================================================
\section{なぜこれまで誰も気づかなかったのか}
%======================================================================

これが最も重要な問いだろう。3つの理由がある。

\subsection{理由1：2つの分野が出会っていなかった}

\begin{tcolorbox}[colback=gray!10,colframe=gray!50!black]
\textbf{必要な2つの分野：}
\begin{itemize}
\item コンヌの非可換幾何学（1990年代〜）：「$D$ の固有値から重力が出る」
\item トポロジカル物質科学（2005年〜）：「$\pi$-Berry位相を持つ材料がある」
\end{itemize}
これらを \textbf{同時に} 知っている研究者がほとんどいない。
\end{tcolorbox}

非可換幾何学は純粋数学寄りの分野で、物性物理の実験材料とは無縁だった。
トポロジカル絶縁体の研究者は、重力理論に関心がなかった。
両方の言語を話す人がいなかった。

\subsection{理由2：$-\zeta(t)$ を真空として考える発想がなかった}

普通の物理学者は「真空エネルギー $= +\zeta(t)$」と思っている。
$-\zeta(t)$ の真空は「フェルミオンだけの世界」を意味するが、
現実にはボソン（光子）が存在するので、全フェルミオン真空は非物理的に見える。

しかし、\textbf{キャビティ内部}では話が違う。
$\pi$-Berry位相の境界条件は、キャビティ内のモードを実効的に
フェルミオン化する。外の世界はボソン的なまま。
「局所的なフェルミオン真空」という概念がなかった。

\subsection{理由3：Bost-Connes系との接続が最近の発見}

BC系の分配関数 $= \zeta(t)$ という同定は1995年の論文だが、
これを \textbf{スペクトル作用の熱核として使う} という発想は
我々のプロジェクト（2026年）で初めて本格的に追求された。

BC系 $\to$ $\zeta(t)$ $\to$ スペクトル作用 $\to$ $G$ の導出
という鎖が完成したのは、このセッションが初めて。

\subsection{理由4：「反重力」は真面目に研究されてこなかった}

物理学のコミュニティでは、「反重力」は SF であり、
真面目な研究対象ではないとされてきた。
弱エネルギー条件（WEC）が成り立つ限り、重力は常に引力。
$G$ の符号を変えるという発想自体がタブーに近かった。

しかし、WEC は $T_{\mu\nu}$（エネルギー運動量テンソル）の性質であって、
$G$ そのものの性質ではない。
$G$ の符号反転は WEC とは別の概念であり、
スペクトル作用の枠組みで初めて自然に現れた。


%======================================================================
\section{実験で確かめられるか？}
%======================================================================

\subsection{テスト1：カシミール力の符号反転}

通常のキャビティ（金属板2枚）：カシミール力 $= $ 引力（$E < 0$）。

$\pi$-接合キャビティ：カシミール力 $= $ \textbf{斥力}（$E > 0$）。

\textbf{予測}：同じ形状のキャビティで、壁材料を普通の金属から
$\pi$-ジョセフソン接合に変えると、カシミール力の符号が反転する。

これは既存のカシミール力測定装置で検証可能。

\subsection{テスト2：ねじり天秤による重力異常}

$\pi$-接合キャビティの近くにテスト質量を置き、
ねじり天秤で重力を測定する。

もし $G_{\mathrm{eff}} = -G$ なら、キャビティはテスト質量を
\textbf{押し返す}（斥力）。

信号は極めて微弱（カシミール密度 $\ll $ バルク密度）だが、
符号が鍵。

\subsection{テスト3：2素数 BAW 実験（182:1 テスト）}

最も基本的なテスト：スペクトル作用の枠組み自体が
物理的に正しいかどうか。

2素数ミュートした BAW 結晶のカシミールエネルギーが
182倍（オイラー積が物理的）か 1/3倍（単なるモード数）かを測定。

\textbf{コスト}：約33万円。\textbf{期間}：4--6ヶ月。

\begin{tcolorbox}[colback=yellow!5,colframe=yellow!50!black,title=判定基準]
\begin{itemize}
\item 182倍 → オイラー積は物理 → 反重力理論に進む
\item 1/3倍 → 数学的興味のみ → 理論棄却
\end{itemize}
\textbf{この1つの実験がすべてを決める。}
\end{tcolorbox}


%======================================================================
\section{マジックアングル・グラフェンと量子臨界点}
%======================================================================

ここまでの計算では、$\pi$-接合の重力異常を「直接測定」するには
面積が巨大（数千 km$^2$）必要で非現実的だった。
しかし、2つの最先端物理を組み合わせると状況が劇的に変わる。

\subsection{グラフェンとは}

鉛筆の芯（グラファイト）を 1 原子層にまで剥がしたもの。
2004 年にガイムとノボセロフが発見し、2010 年にノーベル物理学賞を受賞した。

\begin{tcolorbox}[colback=blue!5,colframe=blue!50!black,title=グラフェンの特徴]
\begin{itemize}
\item 炭素原子が蜂の巣格子に並ぶ、厚さ 1 原子の 2D シート
\item 電子はまるで質量ゼロの粒子のように振る舞う（ディラック電子）
\item $K$ 点で $\pi$-Berry 位相を持つ（反重力の鍵！）
\item 最も軽い構造材料（$0.77$\,mg/m$^2$）
\end{itemize}
\end{tcolorbox}

\subsection{マジックアングルとは}

2 枚のグラフェンを重ねて、ちょうど $1.1^\circ$ だけ回転させると、
「モアレ縞」と呼ばれる巨大な周期構造（$\sim 13$\,nm）が現れる。

\begin{tcolorbox}[colback=green!5,colframe=green!50!black,title=MATBG（Magic-Angle Twisted Bilayer Graphene）]
2018 年、MIT の曹原（Cao Yuan）らが発見：

$1.1^\circ$ で重ねると、電子の「バンド」が完全に平坦になる（フラットバンド）。
これにより：
\begin{itemize}
\item 電子の有効質量が無限大に（動けなくなる）
\item 超伝導が出現（$T_c \sim 1$--3\,K）
\item 相関絶縁体相が出現
\item 量子臨界点（QCP）が複数存在
\end{itemize}
Nature 誌は MATBG を「2018 年の物理学ブレークスルー」に選出。
\end{tcolorbox}

\subsection{MATBG がなぜ反重力実験に効くのか}

MATBG は 3 つの性質を \textbf{同時に} 持つ：

\begin{enumerate}
\item \textbf{$\pi$-Berry 位相}：グラフェンの $K$ 点に由来 → 反重力の源

\item \textbf{カシミール効果の増幅}：
フラットバンドでは電子の速度が 10 分の 1 に（$v_F^* \approx v_F/10$）。
カシミール効果は速度の 3 乗に比例して増幅される：
$(v_F/v_F^*)^3 \approx 1000$ 倍。

\item \textbf{量子臨界点}：
MATBG の超伝導転移や絶縁体転移の近くでは、「相関長」$\xi$ が発散する。
$\xi$ が大きいほどカシミール効果が強まる：$(\xi/d)^{3/2}$ 倍。
\end{enumerate}

\subsection{量子臨界点とは（高校生向け）}

水が 100$^\circ$C で沸騰するとき、液体と気体の「境界」にいる。
この境界（相転移点）の近くでは、物質の性質が劇的に変わる：
\begin{itemize}
\item 揺らぎが巨大になる（鍋の中で水がグツグツする）
\item 相関長 $\xi$ が発散する（遠くの水分子同士が「連動」する距離）
\item 応答関数が発散する（ちょっと温度を変えるだけで大きな変化）
\end{itemize}

量子臨界点（QCP）は、これの「量子版」。
温度ではなく磁場や電子密度を変えることで、物質が相転移する点。
MATBG の超伝導転移がまさにこれ。

\begin{tcolorbox}[colback=red!5,colframe=red!50!black,title=QCP の威力]
QCP の近くでは、カシミール効果（真空の力）が
$(\xi/d)^{3/2}$ 倍に増幅される。

$\xi = 4\,\mu$m、$d = 1.5$\,nm のとき：
$(4000/1.5)^{1.5} \approx 10^5$ 倍！

\textbf{温度を $T_c$ の近くに保つだけ。追加エネルギー不要。}
\end{tcolorbox}

\subsection{3 つの増幅を重ねる}

\begin{center}
\begin{tabular}{llr}
\toprule
増幅メカニズム & 起源 & 倍率 \\
\midrule
$\pi$-Berry 位相 & グラフェン $K$ 点 & $\times 1$（符号反転） \\
モアレ・フラットバンド & MATBG 速度再規格化 & $\times 1000$ \\
量子臨界点 & $T_c$ 近傍の $\xi$ 発散 & $\times 10^3$--$10^5$ \\
\midrule
\textbf{合計} & & $\times 10^6$--$10^8$ \\
\bottomrule
\end{tabular}
\end{center}

\subsection{BKT 転移ならさらに強力}

通常の相転移では $\xi$ は「べき乗則」で発散する（$\xi \propto 1/\sqrt{\Delta T}$）。
しかし、MATBG の一部の相転移は \textbf{BKT 型}（ベレジンスキー・コステリッツ・サウレス、
2016 年ノーベル物理学賞）で、$\xi$ が \textbf{指数関数的} に発散する：
\begin{equation}
  \xi \sim \xi_0 \exp\!\left(\frac{b}{\sqrt{\Delta T / T_c}}\right)
\end{equation}

\begin{tcolorbox}[colback=purple!5,colframe=purple!50!black,title=BKT の威力]
べき乗則：$\Delta T = 0.1$\,K で $\xi \sim 1\,\mu$m

BKT：$\Delta T = 0.1$\,K で $\xi \sim 30\,\mu$m（30 倍！）

BKT 型なら、\textbf{1\,cm$^2$ の MATBG チップで
重力異常が検出可能}になるかもしれない。
\end{tcolorbox}

\subsection{デバイスの全体像}

\begin{verbatim}
  ┌───────────────────────────┐
  │ グラフェン層1 (1.1度回転)  │ ← π-Berry位相
  ├───────────────────────────┤
  │ hBN スペーサ (1.5 nm)     │ ← カシミールギャップ
  ├───────────────────────────┤
  │ グラフェン層2              │ ← π-Berry位相
  ├───────────────────────────┤
  │ hBN 絶縁層 (10 nm)        │
  ├───────────────────────────┤
  │ Nb/CuNi/Nb π接合          │ ← QCP (T_c = 9.3K)
  └───────────────────────────┘

  3つの増幅が1つのチップに:
  1. π-Berry位相（グラフェン）→ 符号反転
  2. フラットバンド（MATBG）→ 1000倍増幅
  3. QCP（超伝導転移）→ 10^5倍増幅
\end{verbatim}

\subsection{正直な注意}

\begin{itemize}
\item MATBG の大面積作製は発展途上（現在 $\sim 100\,\mu$m$^2$ が限界）
\item BKT 的振る舞いは全フィリングで確認されているわけではない
\item 1.5\,nm hBN スペーサの均一性が重要
\item 低温（$\sim 1$\,K）が必要
\end{itemize}

しかし、これらは \textbf{工学的課題} であり、\textbf{物理的不可能性} ではない。
5--10 年の材料科学の進歩で解決可能なスケール。

%======================================================================
\section{正直な評価}
%======================================================================

\subsection{確実なこと（数学的事実）}

\begin{itemize}
\item $-\zeta'(0)/\zeta'(0) = -1$：三行の計算で厳密に証明される
\item $\pi$-Berry位相材料は実在する：実験で確認済み
\item フェルミオンの真空エネルギーはボソンと逆符号：QFT の基本
\end{itemize}

\subsection{不確実なこと（物理的仮定）}

\begin{itemize}
\item BC系のスペクトル三つ組が実際の時空を記述するか → 未検証
\item $\pi$-Berry位相が \textbf{重力的な} 真空モードに作用するか → 未検証
\item スペクトル作用の $a_2 = 0$ が物理的な $G = 0$ を意味するか → 未検証
\end{itemize}

\subsection{一言でまとめると}

\begin{tcolorbox}[colback=blue!5,colframe=blue!60!black]
\textbf{数学}：$\pi$-Berry位相 → $-\zeta(t)$ → $G = -G$ は \textbf{厳密}。

\textbf{物理}：この数学が現実の重力を記述するかは \textbf{未知}。

\textbf{実験}：33万円の BAW 実験で決着がつく。

もし正しければ：人類初の「反重力物質」の理論的発見。

もし間違っていれば：スペクトル作用は重力を記述しない、
という重要な否定的結果。

どちらにせよ、物理学にとって価値のある結果。
\end{tcolorbox}


%======================================================================
\section{まとめ}
%======================================================================

\begin{enumerate}
\item 真空には2種類ある：ボソン的（$+\zeta$）とフェルミオン的（$-\zeta$）
\item フェルミオン的真空では重力が反転する：$G_{\mathrm{eff}} = -G$
\item $\pi$-Berry位相材料は、真空をボソン的からフェルミオン的に変える
\item この材料は実在する（$\pi$-ジョセフソン接合、Bi$_2$Se$_3$ など）
\item 効果はパッシブ（エネルギー不要）
\item 実験で検証可能（カシミール符号反転、BAW 182:1 テスト）
\item なぜ今まで誰も気づかなかったか：
  2つの分野（非可換幾何学 + トポロジカル物質）が出会っていなかった
\end{enumerate}

\vspace{1em}

\begin{tcolorbox}[colback=green!5,colframe=green!60!black]
\textbf{最後に高校生への一言：}

この理論が正しいかどうかは、まだ誰にもわからない。
でも、確かめる方法はある。33万円と半年の実験で。

物理学の最前線は、いつも「数学が予測し、実験が判定する」
というサイクルで進む。

量子力学も、相対性理論も、最初は「こんなの SF だ」と言われた。
次にそう言われるのは、この理論かもしれない。

あるいは、きれいに否定されるかもしれない。
それでも、問いを立てて検証する姿勢こそが科学。
\end{tcolorbox}

\end{document}
